\section[第2章\q 第一课\q 富人不为钱工作]{第2章\\第一课\q 富人不为钱工作}

“爸爸，你能告诉我怎样才能变得富有吗？”

爸爸放下手中的晚报，问道：“你为什么想变富有呢，儿子？”

“因为基米的妈妈会开一辆新凯迪拉克带基米去海滨别墅度周末。基米说要带3个朋友去，但他没有邀请我和迈克，他说这是因为我们是穷孩子。”

“他们真的这么说？”爸爸不相信地问。

“是啊，他们就是这么说的！”我用一种受伤的语调答道。

爸爸默默地摇了摇头，把眼镜往鼻梁上推了推，然后又继续看报纸了。我站在那儿等着答案。

那是1956年，当时我才9岁。由于命运的捉弄，我进了一所公立学校，里面多数学生是富人的孩子。我们这个镇是由最初的甘蔗种植园发展起来的。种植园主和镇上其他有钱人，比如医生、公司老板、银行家，都把孩子送进了这所学校，从一年级一直读到六年级。之后这些孩子通常会被送进私立学校。因为我家住在街的这一边，所以我进了这所学校。如果我家住在街的另一边，我就会去另外一所学校，和那些与我出身差不多的孩子在一起上学。上完六年级之后，我们这些穷孩子会去上公立中学。镇上没有为我们设立的私立中学。

爸爸终于放下了报纸，我敢说他刚才一定是在思考我的话。

“哦，儿子，”他慢慢地开口了，“如果想富有，你就必须学会挣钱。”

“那么怎么挣钱呢？”我问。

“用你的头脑，儿子。”他说着，并微笑了一下，这种微笑意味着“我要告诉你的就这些”，或者“我不知道答案，别为难我了”。

\subsection*{建立合伙关系}

第二天一早，我就把爸爸的话告诉了我最好的朋友迈克。我和迈克可以说是学校里仅有的两个穷孩子。他进这所学校和我一样是由于命运的捉弄。要是有人在学校里划分一条明确的界限，那么我和迈克在和那些有钱的孩子相处时就不会那么局促不安了。其实我们并非真的很穷，但我们感觉很穷，因为其他男孩都有新棒球手套、新自行车，他们的东西都是新的。

爸爸妈妈为我们提供了生活所需的基本用品，像吃的、住的、穿的，但也仅此而已。我爸爸常说：“想要什么东西，自己挣钱买去。”我们想自己挣钱买东西，但确实没有什么工作可以提供给像我们这样的9岁孩子。

“我们怎样才能挣到钱呢？”迈克问。

“我不知道，”我说，“你想跟我合伙吗？”

迈克同意了，于是，在一个星期六的早晨，他成了我生平第一个创业伙伴。整个早上我们都在想怎么挣钱，偶尔我们也会谈起那些“冷酷的家伙”正在基米家的海滨别墅里玩乐。这实在有些伤人，但却是好事，因为它激励我们继续努力去想挣钱的法子。最后，到了下午，一个念头在我们的头脑中闪过，这是迈克从以前读过的一本科普书里得到的启发。我们兴奋地握手，现在我们的合伙终于有了业务内容。

在接下来的几星期里，迈克和我跑遍了街坊四邻，敲开他们的门问他们是否愿意把用过的牙膏皮攒下来给我们。虽然大人们很迷惑，但大多数人还是微笑着答应了。有的人问我们要做什么，我们的回答是：“我们不能告诉您，这是商业秘密。”

几星期过去了，我妈妈变得越来越心烦，因为我们选了靠近洗衣机的地方存放我们的原料。在一个曾用来装番茄酱瓶子的棕色纸盒里，用过的牙膏皮堆得越来越多。

看到邻居们那些脏乱、卷曲的废牙膏皮都到了她这儿，妈妈最后终于采取了行动。“你们两个到底想要干什么？”她问，“我不想再听到‘商业秘密’之类的话，赶快处理掉这些脏东西，否则我就把它们全扔出去！”

迈克和我苦苦哀求，解释说我们已经快攒够了，然后我们就会开始生产。我们告诉她我们正在等一对邻居夫妇用完牙膏，如此一来，我们就可以拿到他们的牙膏皮了。妈妈答应我们延期一周。

我们开始生产的日期提前了。我们的身上承担着压力。我的第一次合作关系，由于货仓收到了妈妈的禁令而受到威胁。迈克的工作变成了告诉邻居们快些用完他们的牙膏，说他们的牙医提倡更频繁地刷牙，我则开始组装生产线。

一天，爸爸载着一个朋友，来看两个9岁的男孩在车道上合力操纵一条全速运转的生产线。到处都是白色的细粉末。在一个长桌上放着几个我们从学校拿回来的装牛奶的纸盒以及家里烤肉用的小炭炉，小炭炉已经被发红的炭烤得闪着红光。

爸爸小心地走过来，由于生产线挡住了通往车库的去路，所以他不得不把车停在路边。当他和他朋友走近时，看到一口钢锅架在炭上，里面的废牙膏皮正在熔化。那时候，牙膏皮不是塑料做的，而是铅制的。

所以一旦牙膏皮上的涂料被烧掉后，铅皮就会熔化，直到变成液体，这时我们就用妈妈的锅垫垫着，将铅液从牛奶盒顶的小孔中倒进去。

牛奶盒里装满了熟石灰，满地的白色粉末就是之前我们将熟石灰和水混合时弄的。由于我一时忙乱，打翻了装熟石灰的袋子，所以弄得整个场地像是被暴风雪袭击了一样。牛奶盒就是熟石灰模的容器。

爸爸和他的朋友注视着我们小心翼翼地把铅液注入到熟石灰模顶部的小孔中。

“小心！”爸爸说。

我顾不上抬头，只是点了点头。

最后，当溶液全部倒入熟石灰模后，我放下钢锅，向爸爸绽开了笑脸。

“你们在干什么？”他带着一丝不解的微笑问道。

“我们正在按照你说的做，我们就要变成富人了！”我说。

“是的，”迈克说，他一边点头一边咧嘴笑着，“我们是合伙人。”

“这些熟石灰模子里是什么东西？”爸爸问。

“看，”我说，“这是已经铸好的一批。”

我用一个小锤子敲击外面的密封物，熟石灰模子被敲成两半，我小心地抽掉熟石灰模的上半部，一个铅制的五分硬币便掉了下来。

“噢，天啊，”爸爸说，“你们用铅造硬币！”

“对啊，”迈克说，“我们正按照你告诉我们的去做。我们正在赚钱。”

爸爸的朋友转过身去纵声大笑，爸爸则微笑着摇着头。除了一堆火和一盒子废牙膏皮，他面前还站着两个灰头土脸的小男孩，他们正咧着嘴笑着。

爸爸让我们放下手里的东西和他坐到屋外的台阶上，然后他微笑着向我们耐心地解释“伪造”一词的含义。

我们的梦想破灭了！“你的意思是说这么做是违法的？”迈克用颤抖的声音问道。

“别去管他们了，”爸爸的朋友说，“这也许是在展现他们的天赋呢。”

我爸爸瞪了他一眼。

“对，这是违法的。”爸爸温和地说，“但是，你们刚才展示了巨大的创造性和新颖的想法，继续努力，我真为你们感到骄傲。”

失望之中，迈克和我呆坐了大约20分钟才开始收拾残局。我们的生意在开始的第一天就结束了。在我清扫熟石灰粉时，我看着迈克，对他说：“我猜基米他们是对的，我们是穷人。”

爸爸正要离开时听到了我的话，“孩子们，”他说，“如果你们放弃了，你们就只能是穷人了。最重要的是你们已经尝试了。大多数人只是夸夸其谈，梦想着发财致富，而你们已经付出了行动。我真为你们骄傲，我要再说一遍，继续努力，不要放弃。”

迈克和我默默地站在那儿，这些话听起来不错，但我们仍然不知道应该做些什么。

“那你为什么不是富人呢，爸爸？”我问。

“因为我选择了当一名老师。老师不该去想怎么发财。我们就是喜欢教书。我希望我能帮你们，但我真的不知道怎么才能赚大钱。”

迈克和我转过身去继续清理现场。

“我了解，”爸爸说，“如果你们想知道如何致富，不要问我，去和你爸爸谈谈，迈克。”

“我爸爸？”迈克苦着一张脸问道。

“对，你爸爸。”爸爸微笑着重复道，“你爸爸和我聘请同一个银行经理，他对你爸爸非常崇拜。他对我提过好几次，说你爸爸在赚钱方面是个天才。”

“我爸爸？”迈克难以置信地问，“那我家为什么没有好车和大房子，就像学校里的那些有钱的孩子家那样呢？”

“拥有好车和大房子不见得就意味着你很富有或你懂得如何赚钱，”我爸爸答道，“基米的爸爸为甘蔗种植园工作。他和我并没有什么不同，他为公司工作而我为政府工作，是公司为他买了那辆车。甘蔗种植园正处于财务困境之中，基米的爸爸过不了多久可能就什么都没有了。但你爸爸不同，迈克。他似乎正在建立一个帝国。我想也许几年之后他就会成为一个非常富有的人。”

听完这番话，我和迈克又兴奋起来了。我们身上充满了干劲，开始清理首次失败的生意所造成的混乱局面。我们一边清理一边制定了一个与迈克的爸爸谈话的计划，例如：该怎样谈，何时谈。问题在于迈克的爸爸工作时间很长，并且经常很晚才回家。他爸爸拥有一个货仓，一家建筑公司，一些连锁商店和3个餐馆。他到很晚才能回家。

在我们清理完之后迈克坐公共汽车回家了。他会在他爸爸晚上回家后和他谈，问他爸爸是否愿意教我们赚钱。迈克答应我，无论多晚，他和他爸爸一谈完就给我电话。

晚上8点30分，电话响了。

“好的，”我说，“下星期六。”我放下了电话，迈克的爸爸同意与我们会面。

星期六早上7点30分，我坐上了公共汽车，向小镇上比较穷困的街区驶去。

\subsection*{课程开始了}

“我每小时付给你10美分。”

即使是以1956年的薪酬标准来看，10美分每小时也是极低的。

我和迈克在那天上午8点与他的爸爸会面了。他已经开始忙碌了，而且在这之前他已经工作了1个多小时。当我走进富爸爸那简朴、窄小而整洁的家时，他的项目监理人刚开着小卡车离开。迈克站在门口迎接我。

“我爸爸正在打电话，他让我们在门廊后面等着。”迈克边开门边说。

当我跨过这座老房子的门槛时，旧木地板发出“吱吱嘎嘎”的响声。

门里面地板上有个破旧的垫子，垫子放在这里是为了隐藏无数脚步经年累月在这块地板上留下的痕迹，它虽然很干净，但还是该换了。

当我走进狭小的客厅时有些害怕，里面塞满了陈旧、发霉而厚重的家具，它们早该成为收藏品了。有两个女人坐在沙发上，她们的年纪比我妈妈大一些，在她们的对面坐着一个穿工作服的男人。他穿着卡其布的衬衫和裤子，衣服烫得很平整，但没有浆过，他手上拿着磨得发光的工作簿。他比我爸爸大10岁左右的样子，我想大概45岁吧。当我和迈克经过他们身边时，他们冲我们微笑，我也有点腼腆地冲他们笑笑。我们向厨房走去，穿过厨房可以到达门廊，在那里可以眺望整个后院。

“他们是什么人？”我问迈克。

“噢，他们是给我爸爸干活的。那个年纪稍大的男人负责管理货仓，那两个女人是餐馆的经理。刚才在门口你看到的是项目监理人，他在离这儿80千米远的一个公路项目中工作。他的另外一个项目监理人正在负责房地产的项目，不过他在你到这儿之前就已经走了。”

“每天都是这样的吗？”我问。

“并不总是，但经常是这样忙的。”迈克笑了笑，拉了一张椅子坐在我身边。

“我问他愿不愿意教我们挣钱。”迈克说。

“哦，那他怎么说？”我急切地问。

“嗯，开始时他露出一种很想笑的表情，然后他说会给我们一个建议。”

“太好了！”我说着，用椅子的两个后腿撑着，把椅子靠着墙翘起来。

迈克也学着我这么做。

“你知道是什么建议吗？”我又问。

“不知道，但很快就知道了。”迈克说。

突然，迈克的爸爸推开那扇摇摇晃晃的纱门走进了门廊，迈克和我跳了起来，不是出于尊敬而是因为被吓了一跳。

“准备好了吗，孩子们？”迈克的爸爸一边问一边拉过椅子坐到我们旁边。

我们点了点头，把椅子扶正在他面前坐下。

他也是个大块头，身高大约有1米80，体重90千克。我爸爸比迈克的爸爸大5岁，他的个子要更高一些，但他们的体重差不多。他们看上去很像同一类人，但气质不同。也许他们的力气都那么大。

“迈克说你们想学赚钱，对吗，罗伯特？”

我赶紧点了点头，但心里有点儿忐忑。在他的微笑和话语的背后似乎隐藏着很强的力量。

“好吧，我说说我的建议：我会教你们，但不像在学校那样。你们为我工作，否则我就不教。因为通过工作我可以更快地教会你们。如果你们只想坐着听讲，就像在学校里一样，那我就是在浪费时间。这就是我的建议，你们可以接受也可以拒绝。”

“嗯……我可以先问个问题吗？”我问。

“不能，你只能告诉我是接受还是拒绝。因为我有太多的事要做，不能浪费时间。如果你不能下定决心，就永远也学不会如何赚钱。机会总是转瞬即逝。知道什么时候要迅速作出决定是一项非常重要的技能。

现在你有一个你想要的机会，但你想进入的这所学校会在10秒钟内开学或者关门。”迈克的爸爸说，脸上带着揶揄的微笑。

“接受。”我说。

“接受。”迈克也说。

“好！”迈克的爸爸说道，“马丁夫人会在10分钟内到这儿。等我和她办完事后，你们就和她坐车去我的小超市，然后就可以开始工作了。

我每小时付给你们10美分，你们每周六工作3个小时。”

“但我今天有一场棒球比赛！”我说。

迈克的爸爸压低声调，用严厉的语气说：“接受或者拒绝。”

“我接受。”我回答道，我决定去工作和学习，不去打棒球了。

\subsection*{30美分以后}

从一个美好的星期六早上9点起，迈克和我正式开始给马丁夫人干活了。马丁夫人是一个慈祥而有耐心的人，她总是说迈克和我使她想起她的两个儿子，她的两个儿子长大后就离开了她。马丁夫人虽然很慈祥，却强调工作应该努力，她让我们不停地干活。她是一个很好的监工，3个小时里，我们把罐装食品从架子上拿下来，用羽毛掸掸去每个罐头上的灰尘，然后重新把它们码好。这工作真的很乏味。

迈克的爸爸，就是我称为“富爸爸”的那一位，拥有9个这样的小型超市，它们是像“7-11”那样的便利店的雏形，当时除了这些小型超市以外附近几乎没有可以买到牛奶、面包、黄油和香烟的杂货店，所以生意还不错。问题是，这是在还没有出现空调的夏威夷，由于炎热，商店不可能关上门。而店的两边有许多停车位，每当一辆车开过或驶进车位，灰尘就漫天扬起飘入店内。

于是，在还没有空调的时代，我们就有事可干了。

此后的3个星期中，每星期六迈克和我向马丁夫人报到并在她那儿工作3小时。中午以前，我们的工作就结束了，她就在我们每人的手中放下3枚硬币。即使是在50年代中期，对于9岁的男孩来说，30美分也并不十分令人激动，因为就算买一本连环画也得花上10美分呢。

第四个星期的星期三，我决定要退出。我答应工作是因为想跟迈克的爸爸学习赚钱，而现在我成了每小时10美分的奴隶。更糟糕的是，自从第一个星期六后我就一直没见到过我们的赚钱老师——迈克的爸爸。

“我要退出。”吃午饭的时候我对迈克说。学校的午饭糟透了，上课也没劲儿，而且我现在几乎一点也不盼着过星期六了。因为对我而言，现在每个星期六换来的仅仅是30美分。

迈克得意地笑了。

“你笑什么？”我既沮丧又气恼。

“我爸说早料到你会退出，他说如果你不想干了就让我带你去见他。”

“什么？”我觉得自己被耍了，于是气愤地问，“他早就在等我去找他？”

“是的，我爸爸可不是一般人，他跟你爸的教育方法不一样。你爸你妈说得多，我爸说得少，不过他早就猜到你会这么说了。你要等到这个星期六，我会告诉他你已经准备好了。”

“你是说我被耍了？”

“不，还不肯定，但有可能。我爸爸会在星期六说明的。”

\subsection*{星期六的排队等候}
我已经准备好要面对迈克的爸爸，就连我的亲爸爸也生气了。我的亲爸爸，我叫他穷爸爸，认为我的富爸爸违反了《童工法》，应该受到调查。

我爸爸要我去争取应有的待遇。每小时至少应该得到25美分。爸爸说如果我得不到加薪，就应该立即辞职。

爸爸气愤地说：“你根本就不需要那份该死的工作。”

星期六早上8点，我又穿过了迈克家那扇摇晃着的纱门。

“坐下等着。”我一进门迈克的爸爸就对我说，说完便转身消失在卧室边的小办公室里。

我环视整个房间，没看见迈克，我感到有些局促，小心地坐到了沙发上，4个星期前我见过的那两个女人笑着给我挪出了点地方。

45分钟过去了，我开始生气了。那两个女人已经见了迈克的爸爸，并且在30分钟之前就离开了。那个年纪大的男人待了20分钟也办完事走了。

一个小时过去了，那天阳光灿烂，我却坐在阴暗、发霉的客厅里，等待着和一个剥削童工的吝啬鬼谈判。我能听见他在办公室里走动、打电话，但他就是不理我。我很想出去，但不知为什么我还是留下来了。

最后，又过了15分钟，正好9点，富爸爸终于走出了他的办公室。

他什么也没说，用手示意我跟着他去那间阴暗的办公室。

“你要求加薪，否则你就不干了？”他边说边把椅子转来转去。

“你不讲信用！”我脱口而出，眼泪差点掉下来。让一个9岁的小男孩去面对一个成年人是会觉得有点害怕。

“你说过如果我为你工作，你就会教我。好，我给你干活，我工作努力，我还放弃了棒球比赛，而你却说话不算数，你什么也没教我！就像镇上人说的，你是一个骗子，你贪心。你就想挣钱，却毫不关心你的雇员。你一点儿也不尊重我，让我等了这么久。我只是个小孩，我应该得到优待！”

富爸爸往转椅里一靠，手摸着下巴盯着我，好像在研究我。

“不错，”他说，“还不到1个月，你已经有点像我的其他雇员了。”

“什么？”我问。我不明白他的话，心里更加委屈了。“我想你会如约教我，你却折磨我。这太残忍了，真的太残忍了！”

“我正在教你。”富爸爸平静地说。

“你教我什么了？什么也没有！”我生气极了，“自从我为那几个小钱干活以来，你甚至都没和我说过话！10美分1小时！哈，我应该到政府那儿告你！”

“你知道，我们有《童工法》，我爸爸可是为政府工作的。”

“哇！”富爸爸叫道，“现在你看上去就像大多数给我干过活的人了，他们要么被我解雇要么辞职不干了。”

“你还有什么可说的？”我说道。作为一个孩子，我觉得自己很有勇气。“你骗了我，我为你工作，而你不守信用，什么都没教我。”

“你怎么知道我什么都没教你？”富爸爸平静地问。

“你从来没和我谈过话，我已经干了3个星期，而你什么也没教给我。”我撅着嘴说。

“教东西一定要说或讲吗？”富爸爸问。

“是呀。”我回答道。

“那是学校教你们的法子，”他笑着说，“但生活可不是这样教你的。我得说生活才是最好的老师。大多数时候，生活不会和你说什么，它只是推着你转，每一次推，它都像是在说：‘喂，醒一醒，我要让你学点东西。’”

“这家伙在说什么呀？”我暗自问自己，“生活推着我转就是生活在对我说话？”现在我知道我必须辞职了，我正在和一个应该被锁起来的家伙说话。

但富爸爸仍在说：“假如你学会了生活这门课程，做任何事情你都会游刃有余。如果你学不会，生活照样会推着你转。人们通常会做两件事，一些人在生活推着他转的同时，抓住生活赐予的每个机会；而另一些人则非常生气，去与生活抗争。他们与老板抗争，与工作抗争，甚至与自己的配偶抗争，他们不知道生活同时也给了他们机会。”

我还是不太明白富爸爸的话。

“生活推着我们所有的人，有些人放弃了，有些人在抗争。少数人学会了这门课程，取得了进步，他们欢迎生活来推动他们，对他们来说，这种推动意味着他们需要并愿意去学习一些东西。他们学习，然后取得进步。但大多数人放弃了，还有一部分人像你一样在抗争。”

富爸爸站起来，关上了那扇嘎吱直响的旧木窗户。“如果你学会了这门课程，你就会成为一个聪明、富有和快乐的人。如果你没有学会，你就只会终生抱怨工作、低报酬和老板，你终其一生希望有个大机会能够把你所有的钱的问题都解决。”

富爸爸抬眼看我是否在听。他的眼光与我的相遇，我们对视着，通过眼神进行着交流。最后，当我接收了他全部的信息后，我将眼睛转开了。我知道他是对的，我责备他，但是是我提出要学习的，我是在抗争。

富爸爸继续说：“如果你是那种没有勇气的人，生活每次推动你，你都会选择放弃。如果你是这种人，你的一生会过得稳稳当当，不做错事、假想着有事情发生时自救，然后慢慢变老，在无聊中死去。你会有许多朋友，他们很喜欢你，因为你真的是一个努力工作的好人。你的一生过得很安稳，处世无误。但事实是，你向生活屈服了，不敢承担风险。你的确想赢，但失败的恐惧超过了成功的兴奋。只有你知道，在你内心深处，你始终认为你不可能赢，所以你选择了稳定。”

我们的眼光又相遇了。我们对视有10秒钟之久，直到相互明白了对方的心意。

“你一直想推动我？”我问。

“可以这样说，但我宁愿说我在让你品尝生活的滋味。”富爸爸笑道。

“什么生活的滋味？”我问，虽然余怒未消，但充满好奇，甚至有点想听他的教诲了。

“你们俩是最先请求我教你们赚钱的人，我有150多个雇员，但没有一个人问过我这个问题。他们只是要求工作和报酬，我从来没有教过他们关于金钱的知识。他们把一生中最好的年华用来挣钱，却不明白到底是为了什么而工作。”

我坐在那儿专心地听着。

“所以当迈克告诉我你们想学赚钱时，我决定设计一个和真实生活相近的课程。虽然我也可以讲得精疲力竭，但你们可能连一个字都听不进去，所以我决定让生活推着你们，这样你们就会记住我的话了，这也就是为什么我每小时只给你们10美分的用意。”

“那么，我又能从每小时10美分的工作中学到什么呢？”我问，“是说你很卑鄙，在剥削工人吗？”

富爸爸向后靠去并开心地笑了起来，然后他说：“你最好改变一下观点，停止责备我，不要认为是我的问题。如果你认为是我的问题，你就会想改变我；如果你认为问题在你那儿，你就会改变自己，学习一些东西让自己变得更聪明。大多数人认为世界上除了自己外，其他人都应该改变。让我告诉你吧，改变自己比改变他人更容易。”

“我不明白。”我说。

“别拿你的问题来责备我。”富爸爸说，他开始有些不耐烦了。

“可你每小时只给我10美分。”

“那么你学到了什么？”他笑着问。

“你很卑鄙。”我顽皮地笑了笑。

“瞧，你还是觉得问题在我这儿。”富爸爸说。

“可的确是这样呀。”

“好吧，如果继续这种态度，你就什么都学不到。如果你仍认为问题在我这儿，你该怎么办？”

“嗯，如果你不提高我的工资，更尊重我并教我赚钱，我就辞职。”

“说得好，”富爸爸说，“大部分人会这么干，他们辞职，然后去找另一份工作，期望得到更好的机会、更高的报酬，他们认为一份新的工作或更高的报酬会解决问题。在大多数情况下，这是不可能的。”

“那什么能解决问题呢？”我问，“接受这可怜兮兮的每小时10美分还要报以微笑吗？”

富爸爸笑了。“有些人会这么做，只因为他们和他们的家庭需要钱而接受这份工资。但他们所做的也只是等待，等待加薪，因为他们认为更多的钱能解决问题。大部分人接受这样的工资，还有一些人会再找一份工作，仍旧干得很努力，但仍只能得到很少的报酬。”

我坐在那里，眼睛盯着地板，开始理解富爸爸给我们上的这一课。

我感到这的确是生活的原味。最后，我抬起头，又重复了前面的问题：“那么用什么来解决问题呢？”

“用这个，”他说着轻轻地拍了拍我的脑袋，“用你两个耳朵之间的家伙。”

就在那一刻富爸爸和我们分享了使他区别于他的职员和穷爸爸的最关键的东西——这使他最终成为夏威夷最富有的人之一。而我受过良好教育的爸爸则一生都在与财务问题抗争。富爸爸非凡的观念使他的一生都与众不同。

富爸爸不厌其烦地讲述这个观点，一遍又一遍，这就是我称之为“第一课”的内容。

穷人和中产阶级为钱而工作。

\subsection*{穷人和中产阶级为钱工作}

在那个阳光明媚的星期六上午，我学习了一种与穷爸爸教给我的完全不同的观念。在我9岁的时候，我意识到两位爸爸都希望我去学习，鼓励我去研究，但研究的内容不同。

我那受过高等教育的爸爸建议我像他那样做。“儿子，我希望你努力学习，取得好成绩，这样你就能在大公司里找到一份稳定的工作，而且会收入不菲。”富爸爸却希望我去研究钱的运动规律，好让钱为我所用。在他的指导下，我会在生活中而不是在教室里学习这些课程。

富爸爸继续给我上第一课：“我很高兴你为每小时10美分的报酬生气，如果你不生气而是高兴地接受了，那我只能说我没法教你。你看，真正的学习需要精力、激情和热切的愿望。愤怒是其中一个重要的组成部分，因为激情正是愤怒和热爱的结合体。说到钱，大多数人都希望稳稳妥妥地挣钱，这样他们才感到安全。关于钱，他们没有激情，有的只是恐惧。”

“这就是他们接受低工资的原因？”我问。

“是呀，”富爸爸说，“有人说我剥削工人，因为我比甘蔗种植园和政府付给员工的薪水少。我说是他们自己剥削自己，罪魁祸首是他们的恐惧，而不是我。”

“但你不觉得你应该多给他们一点儿薪水吗？”我问。

“没这必要。而且，钱多了也解决不了问题。比如你爸爸，挣钱也不少，但仍会欠债。对大多数人而言，给他们的钱越多，他们欠的债也就越多。”

“这就是每小时10美分的原因，”我笑了，“课程的一部分？”

“没错。”富爸爸也笑了，“你瞧，你爸爸上了大学而且受到很好的教育，所以他有希望得到一份高薪的工作。他的确做到了，但他还是为钱所困，原因就是他在学校里从来没学过关于钱的知识。而且最大的问题是，他相信工作就是为了钱。”

“你不这么认为吗？”我问。

“不，当然不是，”富爸爸回答，“如果你想为钱而工作，那就待在学校里吧，那可是一个学习这种事的好地方。但是如果你想学习怎样让钱为你工作，那就让我来教你。不过首先你得想学。”

“难道不是每个人都想学吗？”我问。

“不是，”他说，“原因很简单，学习为钱工作很容易，特别是当你一谈到钱就觉得恐惧时，学习为钱工作就更容易了。”

“我不明白。”我皱着眉头说。

“别担心，你只需知道，正是出于恐惧的心理，人们才想找一份安稳的工作。这些恐惧有：害怕付不起账单，害怕被解雇，害怕没有足够的钱，害怕重新开始。为了寻求保障，他们会学习某种专业，或是做生意，拼命为钱而工作。大多数人成了钱的奴隶，然后就把怒气发泄在他们老板身上。”

“学习让钱为我工作和上面这种不一样吗？”我问。

“当然了，”他说，“绝对不一样。”

在夏威夷这个美丽的星期六早晨，我们静静地坐着。我的朋友们应该已经开始新一季的棒球联赛了，但不知为什么，我突然开始庆幸自己决定干这份每小时10美分的工作了，我感到我会学到我的朋友们在学校里学不到的东西。

“准备好了吗？”富爸爸问。

“当然。”我咧开嘴笑了。

“我已经遵守了诺言，带你去看了你未来的生活。”富爸爸说，“你现在才9岁时，就已经有了为钱工作的体验了。你只需把上个月的生活重复50年，就会知道大多数人是如何度过一生的了。”

“我不太懂。”我说。

“你两次等着见我时有什么感觉？上次是被雇用，这次是要求加薪。”

“真可怕。”我说。

“如果人们选择为钱工作，这就是他们将要过的生活。”

“每次结束3小时的工作，马丁太太给你3个硬币，这时你又有什么感觉？”

“我觉得钱不够。感觉就像什么也没得到似的，真让人失望。”

“这也正是大多数雇员拿到工资单时的感觉，他们还要扣掉税和其他一些支出。至少，你拿到的还是100％的工资。”

“你是说工人们拿到的不是全部的工资？”我吃惊地问。

“当然不是，”富爸爸说，“政府要先拿走属于它的那份。”

“它怎么拿呢？”

“通过税收，”富爸爸说，“你挣钱时得缴税，花钱时也得缴税。你存钱时得缴税，你死时还得缴税。”

“政府怎么能这样？”

“富人就不会这样，”富爸爸微笑着说，“只有穷人和中产阶级是这样。我敢打赌我赚的比你爸爸多，但他缴的税比我多。”

“怎么可能呢？”我问道。作为一个9岁的男孩，是不会理解这些的。

“为什么有的人会让政府这么对待自己呢？”

富爸爸坐在那儿沉默不语，我猜他希望我认真地听而不是说一些无聊的话。

于是我安静下来。我不喜欢刚才听到的事情。我知道爸爸总是抱怨税收太高了，但没有采取任何措施改变这种局面。生活是否也在推着他？

富爸爸缓慢而沉默地摇着座椅，看着我。

“真的准备好跟我学习了吗？”他问。

我郑重地点了点头。

“就像我说的，这里头有不少东西要学。学习怎样让钱为你工作是一个持续终生的过程。大多数人上了4年大学后，学习也就到头了。可我知道我会用一辈子去研究钱这东西，因为我研究得越深，就越发现我还有更多的东西要学习。大多数人从不研究这个问题，他们去上班，挣工资，平衡收支，仅此而已，他们不明白自己为什么老缺钱，于是以为多挣点钱就能解决问题，但几乎没有人意识到缺乏财商教育才是问题的关键。”

“那我爸爸总为税头疼也是因为他不懂钱吗？”我疑惑地问。

“税只是学习如何让钱为你工作的一个极小的部分。今天，我只想弄清楚你是否仍有热情去了解钱这东西。大多数人都没有这样的愿望，他们只想进学校，学习一门专业技能，然后轻松工作、挣大钱。到他们某一天醒来，发现已面临严重的财务问题时，他们已经不能停止工作了。这就是只知道为钱工作而不学习如何让钱为自己工作的代价。现在你还有热情学习吗？”富爸爸问。

我点了点头。

“好，”他说，“现在回去干活，这次我一分钱也不会给你。”

“什么？”我大吃一惊。

“听着。一分钱也不给。每星期六你照样干3个小时，但不会再有每小时10美分的报酬了。你说你想学习不为钱工作，所以我什么都不给你。”

我几乎不相信自己的耳朵。

“我已经和迈克谈过了，他已经开始免费干了，掸干净罐头上的尘土再把它们重新码好。你最好快点去和他一块儿干。”

“这不公平，”我叫道，“你总得给点什么呀。”

“你说过你想学习。如果你现在不学，将来就会像坐在客厅里的那两个女人和那个男人一样，为钱工作并且希望我别解雇他们。或是像你爸爸那样，挣了很多钱却眼看着债台高筑而毫无办法，希望靠拥有更多的钱来解决问题。如果你想这样，我可以按照原来的约定每小时付给你10美分。你也可以像大人那样，抱怨这里工资太低，辞职另找工作。”

“那我该怎么做呢？”我问。

富爸爸拍了拍我的头，“动动脑子，”他说，“如果你好好想一想，就会感谢我给你这个机会，让你成为有钱人。”

我站在那儿，依旧不相信自己已经答应了这个不公平的交易。我是来要求加薪的，现在却被告知以后要白干。

富爸爸又拍了拍我的头，说：“慢慢想去吧，现在回去开始工作。”

\subsection*{富人不为钱工作}

我没对我爸爸说我没工钱了，他是不会理解的，而且我也不想向他解释连我自己也没完全明白的事。

在接下来的3个星期里，我和迈克每个星期六白干3小时。这工作不再让我心烦，过程也容易些了。只是无法参加棒球赛以及不能再买连环画让我耿耿于怀。

富爸爸在第三个周末的中午来了。我们听见他的卡车进了停车场，然后发动机熄火了。他走进店里与马丁太太拥抱致意。他看了店面的销售情况后，走向冰激凌柜，取出两个冰激凌，付了钱，然后向我和迈克打了个手势。

“孩子们，我们出去走走。”

闪开来来往往的汽车，我们穿过街道，又走过一大片草地，草地上有许多大人正在打垒球。最后我们坐到草地深处的一张野餐桌前，富爸爸把冰激凌递给我和迈克。

“还好吗？”他问。

“挺好的。”迈克说。

我也点了点头。

“那学到了什么没有？”

迈克和我面面相觑，一起耸了耸肩。

\subsection*{避开一生中最大的陷阱}
“好吧，孩子们，你们最好开始开动脑筋。你们正在学习一生中最重要的一课。如果学好了这一课，你们将永享自由和安宁；如果没有学好，你们就会像马丁太太和在草地上打垒球的大多数人一样度过一生。

他们为了一点点钱而勤奋工作，深信有工作就有了保障，盼着一年3个星期的假期和工作45年后才能获得的一小笔养老金。如果你们喜欢这样，我就把工资提到每小时25美分。”

“但他们都是努力工作的好人啊，你在嘲笑他们吗？”我问道。

一丝笑容浮上了富爸爸的脸庞。

“马丁太太对我就像妈妈一样，我决不会那么残忍地看待她。我上面的话可能听起来很无情，那是因为我在尽力向你们说明一些事情。我想拓宽你们的视野，让你们看清一些东西。这些东西甚至连大多数成年人也从没认清，因为他们眼界太狭窄了。大多数人从未认识到他们是在陷阱之中。”

迈克和我还是不太明白他的意思。他的话听起来很无情，然而我们能感到他确实是急于想让我们明白。

富爸爸又笑着说：“25美分每小时怎么样？这样是否让你们心跳加速？”

我摇摇头说：“不会啊。”可事实上，25美分每小时对我而言可真是不小的数目啊！

“那么，我每小时给你1美元。”富爸爸说，脸上露出狡黠的笑容。

我的心开始狂跳，脑袋里有个声音在喊：“接受，快接受。”但我不相信我所听到的，所以什么也没说。

“好吧，每小时2美元。”

我这个9岁孩子的大脑和心脏几乎要爆炸了。毕竟这是1956年，每小时2美元的薪水将使我成为世界上最有钱的孩子！我无法想象能挣到这么多钱。我想说“好的”，真想和他成交。我似乎已经看见一辆新自行车，一副新棒球手套，以及当我亮出钞票时同学们羡慕的表情。最重要的是，基米和他的朋友再也不能叫我穷人了，但不知什么原因我仍然没有开口。

也许我的脑袋已经热昏了，但在内心深处，我非常想要那每小时的2美元。

冰激凌化了，顺着我的手流了下来。地上留下一滩黏黏的香草和巧克力味的冰激凌，蚂蚁正在享受它们。富爸爸看到两个孩子盯着他，眼睛睁得大大的，脑子里却空空如也。他是在考验我们，而且他也知道我们很想接受这笔交易。他知道每个人的灵魂都有软弱、贫乏的一面，也有强大、坚定、无法被金钱收买的一面。问题在于哪一部分更强大。他在一生中考验了成百上千的人，每一次招聘面试就是一次考验。

“好，5美元每小时。”

我的内心突然平静了，想法发生了转变。这个出价太高了，高得有些离谱。在1956年，就算成年人也没有几个能每小时挣5美元的。诱惑突然不见了，我恢复了平静。我轻轻转向左边去看迈克，他也在看我。

我灵魂中软弱而贫乏的一面沉默了，无法用钱收买的一面占了上风。面对钱，我开始心安神定。我知道迈克也一样。

“很好，”富爸爸轻声说，“大多数人都希望有一份工资收入，因为他们都有恐惧和贪婪之心。一开始，没钱的恐惧会促使他们努力工作，得到报酬后，贪婪或欲望又让他们想拥有所有用钱能买到的好东西。于是就形成了一种模式。”

“什么模式？”我问。

“起床，上班，付账，再起床，再上班，再付账……他们的生活从此被这两种感觉所控制：恐惧和贪婪。给他们更多的钱，他们就会以更高的开支重复这种循环。这就是我所说的‘老鼠赛跑’。”

“就没有另外一种模式吗？”迈克问。

“有，”富爸爸缓缓说道，“但只有少数人知道。”

“到底是什么模式？”迈克问。

“这就是我希望你们能在工作和跟我学习的过程中找到的东西。也是我不给你们任何工资的原因。”

“有什么建议吗？”迈克问，“我们厌倦了辛苦地工作，尤其是什么报酬都没有。”

“哦，第一步是讲真话。”富爸爸说。

“我们可没撒谎。”我说。

“我没说你们撒谎，我是说要弄清真相。”富爸爸反驳道。

“关于什么的真相？”我问。

“你真正的感觉，”富爸爸说，“你无需告诉别人你的感觉，只有你自己知道。”

“你是说这公园里的人，还有那些为你工作的人，像马丁夫人，他们都没弄清楚自己的感觉？”我问。

“我想是的。他们害怕没有钱，更没有直面这种恐惧，对此他们虽然在情感上有所反应但并没有动脑筋想办法。”富爸爸说着拍拍我们的头，“他们手中有点小钱，可享乐、欲望和贪婪会立刻控制他们，他们会再次作出反应，仍然是不假思索。”

“所以，他们的感情代替了他们的思想。”迈克说。

“正是如此，”富爸爸说，“他们并不清楚自己真正的感觉，只是作出反应，而不去思考。他们感到恐惧，于是就去工作，希望钱能消除恐惧，但没有奏效。于是，恐惧追逐着他们，他们只好又去工作，再一次期望钱能平复这种恐惧，但还是没有成功。恐惧使他们落入工作的陷阱，挣钱——工作——挣钱，希望恐惧就此烟消云散。但每天他们起床时，就会发现恐惧又与他们一起醒来了。恐惧使成千上万的人彻夜难眠，忧心忡忡。所以他们又起床去工作了，希望薪水能消除噬咬他们灵魂的恐惧。钱主宰着他们的生活，他们拒绝去分辨真相，钱控制了他们的情感和灵魂。”

然后富爸爸静静地坐着，让我们自己理解他的话。迈克和我听到了他的话，但不能完全明白他的意思。我只知道我很奇怪大人们为什么总是急急忙忙去工作，而工作看起来并没什么乐趣可言，而且他们也不快活，但好像总有些东西逼着他们去工作。

看到我们已经尽力地理解了他的话后，富爸爸说：“我希望你们能避开这个陷阱，这就是我真正想教你们的，而不只是发财，因为发财并不能解决问题。”

“不能吗？”我惊奇地问。

“不能。现在让我谈谈另一种感情：欲望。有人把它称为贪婪，但我更喜欢用欲望这个词。希望拥有一些更好、更漂亮、更有趣或更令人激动的东西，这是相当正常的。所以人们也为了实现欲望而工作。他们认为钱能买来快乐，可用钱买来的快乐往往是短暂的，所以不久他们就需要更多的钱来买更多的快乐、更多的开心、更多的舒适和更多的安全感。于是他们继续工作，以为钱能安抚他们备受恐惧和欲望折磨的灵魂，但实际上钱是无法做到这一点的。”

“即使是富人也这样吗？”迈克问。

“富人也是如此。”富爸爸说，“事实上，许多人致富并非出于欲望而是由于恐惧，他们认为钱能消除贫困带来的恐惧，所以他们积攒了很多的钱，却发现恐惧感更加强烈了。他们又开始害怕失去钱。我有一些朋友，他们已经很有钱了，但还在拼命工作。我还认识一些百万富翁，他们现在甚至比他们穷困时还要恐惧，他们害怕失去所有的钱。他们越富有，这种感觉就越强烈。他们灵魂中软弱贫乏的一面总是在大声尖叫，他们不想失去大房子、车子和钱带给他们的上等生活。他们甚至担心一旦没钱了，朋友们会看不起他们。许多人变得绝望而神经质，尽管他们很富有。”

“那穷人是不是要快活一点？”我问。

“我可不这么认为，”富爸爸回答说，“不谈钱就像依赖钱一样是一种精神上的疾病。”

就在这时，镇上的乞丐经过我们的桌子，在一个大垃圾箱里翻找起来。我们3个人极有兴趣地注视着他，刚才我们几乎没意识到他在这里。

富爸爸从钱包里掏出1美元，向那个老乞丐招了招手。乞丐看到钱，立即走过来接钱，他对富爸爸千恩万谢，然后就欣喜若狂地跑了。

“他与我的大多数雇员并没有太大差别，”富爸爸说，“我遇到过很多人，他们说‘我对钱没兴趣’，可他们却每天工作8小时。这只能说明他们并没有说出真相，如果他们对钱没兴趣，又何必工作呢？这种人比敛财的人病得更重。”

我坐在那儿听着富爸爸的话，脑中无数次地闪现我的亲爸爸曾说过的话：“我对钱不感兴趣。”他常常这么说。他还经常说：“我工作是因为我热爱这个职业。”用这句话来掩藏他内心真实的感受。

“那我们该怎么办呢？”我问，“不为钱工作直到不再有恐惧和贪婪吗？”

“不，那只会浪费时间。正是因为有感情，我们才成为人。感情使我们更加真实，它是我们行动的动力。忠实于你的感情，以你喜欢的方式运用你的头脑和感情，不要让它们控制你。”

“噢！”迈克叫了起来。

“不要为我的话忧心忡忡，将来你们会知道它很有道理。好好观察你的感情，别急于行动。大多数人并不知道是他们的感情代替了他们进行思考，感情只是感情，你还必须学会抛开感情来思考。”

“你能给我们举个例子吗？”我问。

“可以。”富爸爸回答说，“当一个人说‘我得去找份工作’，这就很可能是他的感情代替他在思考。是害怕没钱的感觉催生了找工作的想法。”

“但是人们确实需要钱来付账呀。”我说。

“的确如此，”富爸爸笑着说，“所以，我说感情常常过多地代替了思考。”

“我不懂。”迈克说。

“比如说吧，”富爸爸说，“如果害怕没钱花，也先不要去找工作，挣几个小钱去消除恐惧，而要先问问自己：一份工作是最终消除这种恐惧的最佳解决办法吗？依我看，答案是‘不是’，从人的一生来看更是如此。工作只是试图用暂时的办法来解决长期的问题。”

“但我爸爸总是说‘去上学，取得好成绩，这样你就能找到一份安稳的工作’。”我有些迷惑地说。

“是啊，我懂他的意思。”富爸爸笑着说，“大多数人都这么给别人建议，而且对于大多数人来说这也确实是个好主意。但人们仍是基于恐惧才给出这样的建议的。”

“你是说我爸爸这么说是因为害怕？”

“是的，他担心你将来挣不到钱，在这个社会上过得不好。别误解我的话，他爱你而且希望你能够一帆风顺。我认为他的担心不无道理。

教育和工作是很重要的，可它们对付不了恐惧。实际上，促使他每天去上班挣钱的恐惧也使得他热衷于让你去上学。”

“那你有什么建议？”我问。

“我想教你们支配钱，而不是害怕它，这是在学校里学不到的。如果不学，你就会变成金钱的奴隶。”

这听起来很有道理，他想扩展我们的视野，让我们看到马丁太太、他的雇员和我爸爸都看不到的东西。富爸爸用了似乎很无情的例子，但这些例子让我终生难忘。在那一天我的视野大开，开始注意到大多数人所面临的“陷阱”。

“你看，我们在根本上都是雇员，只是层次不同而已。”富爸爸说，“我只希望你们有机会避开由恐惧和欲望组成的陷阱，按照你们喜欢的方式利用恐惧和欲望，而不要让它们控制你们。这就是我想教你们的。我对教你们挣大钱不感兴趣，那解决不了问题。如果你们不先控制恐惧和欲望，即使你们获得高薪，也只不过是金钱的奴隶而已。”

“那我们怎么才能避开陷阱呢？”

“造成贫困和财务问题的主要原因是恐惧和无知，而不是经济环境、政府或者富人。人们自身的恐惧和无知使他们困在陷阱里，所以你们应该去上学、接受高等教育，而让我来教你们怎样不落入陷阱。”

谜底渐渐揭晓。我爸爸受过高等教育，事业有成，但学校从来没有告诉他如何处理金钱和恐惧。很显然，我可以从两个爸爸那里学到内容不同但同样都很重要的东西。

“你刚才讲的是对于没钱的恐惧，那么，对钱的欲望又会怎样影响我们的想法呢？”迈克问。

“在我用更高的工资诱惑你们时，你们感觉怎样？非常想要吗？”

我们点了点头。

“但你们没有屈服于自己的感觉，你们没有立刻作出决定。这一点最重要。我们总是会有恐惧、贪婪的时候。从现在开始，对你们来说最重要的是，运用感情作长远打算，别让感情控制了思想。大多数人让恐惧和贪婪来支配自己，这是无知的开始。因为恐惧和贪婪，大多数人一生都在追求工资、加薪和职业保障，从来不问这种感情支配思想的生活之路将通向哪里。就像一幅画表现的：驴子拉车，因为主人在它面前挂了个胡萝卜。主人清楚自己想要去哪里，而驴子却只是在追逐一个幻影。但第二天驴子依旧会去拉车，因为又有胡萝卜放在它的面前。”

“你的意思是，当我幻想新棒球手套、糖果和玩具时，就像那头驴子和它面前的胡萝卜一样？”

“不错。当你长大后，你想要的玩具会更贵，会变成要让你的朋友羡慕的汽车、游艇和大房子，”富爸爸笑着说，“恐惧把你推出门外，欲望又开始召唤你，诱惑你去触礁。这就是陷阱。”

“那我们该怎么做呢？”迈克问道。

“无知使恐惧和欲望更加强烈，这就是为什么很多有钱人越有钱就越害怕。钱就是驴子面前的胡萝卜，是幻象。如果驴子能了解到全部事实，它可能会重新想想是否还要去追求胡萝卜。”

富爸爸继续解释说人生实际上是在无知和觉醒之间的一场斗争。

他说一个人一旦停止了解有关自己的知识和信息，就会变得无知。

这种斗争实际上就是你时刻都要做的一种决定：是通过不断学习打开自己的心扉，还是封闭自己的头脑。

“看，学校非常重要。你去学校学习一种技术或一门专业，并成为对社会有益的人。每一种社会文明都需要教师、医生、工程师、艺术家、厨师、商人、警察、消防员、士兵。学校培养了这些人才，所以我们的社会才能蒸蒸日上。”富爸爸说，“但不幸的是，对许多人来说，离开学校是学习的终点而不是起点。”

然后是长长的沉默。富爸爸依旧微笑着，我仍然没弄明白他说的那些话。但就像其他伟大的老师一样，他们的话多年来会一直给人们以启发，即使他们去世了那些话的作用也不会减小，直到现在我还在回味富爸爸话中的道理。

“今天我有点无情，”富爸爸说，“但无情是有原因的，我希望你们永远记住这次谈话，希望你们多想想马丁太太，多想想那头驴子。永远不要忘记，你有两种感情——恐惧和欲望，如果你让它们来控制你的思想，你就会落入一生中最大的陷阱。一直生活在恐惧中，从不追求自己的梦想，这是残酷的。为钱拼命工作，以为钱能买来快乐，这也是残酷的。半夜醒来想着还有许多账单要付是一种可怕的生活方式，以工资的多少来决定过什么样的生活不是真正的生活。认为工作会给你带来安全感其实是在欺骗自己。这些都很残酷，但我希望你们能尽可能地避开这些陷阱。我看过钱如何控制人们的生活，别让这些问题发生在你们身上，别让钱支配你们的生活。”

一个垒球滚到了桌子底下，富爸爸把它捡起来扔了回去。

“无知与恐惧和贪婪有什么关系？”我问。

“对钱的无知导致了恐惧和贪婪，”富爸爸说，“我可以给你们举一些例子。一个医生，想多挣些钱让家人生活得更好，于是就提高了收费，这就导致人们医疗保健费用的增加。这极大地损害了穷人的利益，所以穷人的健康状况要比富人差。

“由于医生提高收费，律师也提高了收费；由于律师提高收费，学校的老师也想增加收入，这就使政府提高了税收。这种循环不断继续下去，不久，在富人和穷人之间就有了一条可怕的鸿沟，不安定就会出现。当鸿沟大到极点时，一个社会就会崩溃。美国现在的情况就是这样，这种历史一再重演，是因为人们没有以史为鉴。我们只记住了历史事件发生的时间和名称，却没有记住教训。”

“涨价一定不对吗？”我问。

“在一个民众受到良好教育和政府治理有方的社会中不会涨价，实际上应该降价。当然，通常在理论上是如此。价格上涨的原因是由无知引起的贪婪和恐惧。如果学校教学生关于钱的知识，社会就有可能更富足，物价也会更低廉。但学校只关注教学生如何为钱工作，而不是掌握金钱。”

“但我们不是有商学院吗？”迈克问，“你不是还鼓励我进商学院读MBA吗？”

“是的，”富爸爸说，“但不管怎么说，商学院是专门训练那些精于计算的人的，可千万不要让他们做生意。他们只会看看数字，解雇员工并把生意搞糟。我知道这一点是因为我也雇用了这样的人。他们所想的只是降低成本、提高价格，这只会带来更多的问题。精于计算是重要的，我希望更多的人能够懂得它，但这并不是全部。”富爸爸生气地补充道。

“那该怎么办呢？”迈克问。

“要学会让感情跟随你的思想，而不要让思想跟随你的感情。当你们控制了感情，同意免费干活时，我就知道你们还有希望。当我用更多的钱诱惑你们时，你们又一次控制住了感情，你们在学习用头脑思考而不是任由感情控制你们。这是第一步。”

“这一步为什么很重要？”我问。

“噢，这就要由你们自己来找答案了。如果你们想学，我就会把你们带上这条布满荆棘的道路，大多数人都会选择避开它。我会带你们去大多数人都害怕去的地方，跟着我，你们就要抛弃为钱工作的观念，学习如何让钱为你们工作。”

“如果我们跟着你结果会怎样？如果我们同意跟你学，又能学到什么呢？”我问。


“你们会得到《柏油孩子》\footnote{《柏油孩子》讲的是：一个农夫种的庄稼经常被野兔偷吃，所以他用柏油做了一个孩子像立在地头。野兔看到很好奇，上前逗弄柏油孩子却被粘住了，正好被农夫抓住。野兔情急之中跟农夫说，惩治它最好的办法是将它扔进荆棘丛，因为那样会扎得它痛不欲生。农夫信以为真，将野兔扔进了荆棘丛，然后野兔就逃走了，因为荆棘丛就是它的家。}里那只野兔得到的东西——自由。”

“那会是一条布满荆棘的路吗？”

“是的，”富爸爸说，“所谓的荆棘之路就是指我们的恐惧和贪婪。

走进我们的恐惧，直面我们的贪婪、弱点和缺陷是唯一的出路。这条路需要你用心去确定你的思想。”

“确定思想？”迈克不解地问。

“是的，确定我们该怎样思考而不只是对感情作出反应。不要因为害怕付不起账单就起床去工作，希望以这种方法来解决问题。你要花时间去思考这个问题：更努力地工作是解决问题的最好方法吗？大多数人都害怕知道真相——他们被恐惧所支配——不敢去思考，就出门去找工作了。他们被柏油孩子‘粘’住了。这就是我说的确定你的思想。”

“我们怎样才能做到呢？”迈克问。

“那是我将来要教你们的。我会教你们选择一种思想，而不是条件反射式地去行动，就像匆忙喝完早餐咖啡就跑出去工作一样。

“记住我之前说过的话：工作只是面对长期问题的一种暂时的解决办法。大多数人心里只有一个问题，并且亟待解决，那就是月末要付账了，账单就像那个柏油孩子。钱控制了他们的生活，或者确切地说是，对钱的无知和恐惧控制了他们的生活。所以他们就像他们的父母一样，每天起床之后就去工作挣钱，没有时间问问自己：还有什么别的法子吗？他们的思想被他们的感情，而不是他们的头脑控制着。”

“你能说说感情控制思想和头脑控制思想的区别吗？”迈克问。

“噢，当然。我总是听到这种话，”富爸爸说，“有人会说‘人人都必须去工作’、‘富人是骗子’，或是‘我要换份工作，我应该得到更高的工资，你不能任意摆布我’，还有‘我喜欢这份工作，因为它很稳定’，而不是说，‘我失去了什么’，只有这样的话才能让你避免感情用事，并留给你仔细思考的时间。”

我得承认，这的确是重要的一课，即知道人什么时候是在抒发感受，什么时候是在清楚地表达思想。这一课让我终生受益，尤其是每当我的话也是出于情感上的反应而非出于深思的时候。

在我们往回走的路上，富爸爸解释说富人实际上是在“造钱”，他们从不为钱而工作。他接着说当我和迈克用铅铸5分钱的硬币时，我们是想“造钱”，实际上我们的想法和富人的想法很接近。问题是我们这么做不合法，只有政府和银行做这种事才是合法的，我们不行。他解释说有些挣钱的方式是合法的，有些是非法的。

富爸爸继续说，富人知道钱是虚幻的东西，就像挂在驴子面前的胡萝卜一样。正是恐惧和贪婪使无数人紧抓着这个幻觉不放，还以为它是真实的。钱的确是造出来的，正是由于对这种幻觉的信任以及人们的无知才使人们作出许多经不起推敲的计划。“事实上，”富爸爸说，“从许多方面来说，驴子的胡萝卜比钱有价值。”

他说到美国当时实行的是金本位，每一张美元都对应着一定重量的黄金。他对美国将要取消金本位以及美元将不再与黄金挂钩的传言十分感兴趣。

“如果真的发生这种事，孩子们，地狱之门就快开了。穷人、中产阶级和无知的人的生活将被摧毁，因为他们相信钱是真的财富，相信他们效力的公司或政府会保障他们的生活。”

我们没有完全弄明白这席话的含义，但多年以后，富爸爸的话在越来越多的地方应验了。

\subsection*{看见了别人看不见的}
富爸爸上了停在小店外的小卡车，这时他对我们说：“继续工作，孩子们，你们越快忘记你们的工资，你们未来的生活就会越轻松，继续用你们的头脑思考，不求回报地工作，很快就会发现比拿工资更挣钱的方法。你们会看到别人看不见的东西。机会就摆在人们面前，但大多数人从来看不到这些机会，因为他们忙着追求金钱和安定，所以只能得到这些。如果你们能看到一个机会，就注定你们会在一生中不断地发现机会。那时，我会教你们其他的事。学会了这些，你们就能避开生活中最大的陷阱，永远不会被那个柏油孩子粘住了。”

我和迈克收拾好东西与马丁太太道了别。我们走回公园，又坐到那张长椅上，花了几个小时考虑和讨论富爸爸的话。

在接下来的一个星期里，我和迈克一直在思考和讨论这些问题。又过了两个星期，我们仍旧在思考、讨论，同时免费为富爸爸工作。

在第二个星期六工作结束时，我向马丁太太道别，依依不舍地看着架子上的连环画。每星期六没有了30美分的收入，我就没有钱去买连环画了。而就在马丁太太和我们说再见时，她做了一件我以前从未见她做过的事，我的意思是我以前也见她这样做过，只是我没太在意。

马丁太太把连环画的封面撕成两半，然后把上面一半留下，下面一半扔进了棕色的书橱。我问她这是做什么，她说：“我要把这些没有卖掉的旧书扔掉。等书商送新书来，我会把这半边封面交给他，作为没有卖掉的证明。他一小时后就到。”

我和迈克等了一个小时，书商终于来了。我问他能否把那些他不要的连环画送给我们。他回答：“如果你们在这家店干活，而且保证不把它们卖掉，我就送给你们。”

于是，我们成交了。迈克家的地下室里有个空房间，我们把它清理出来，把几百本连环画搬了进去。很快我们的连环画阅览室就开始营业了。我们雇了迈克的妹妹——她很爱读书——来做图书管理员。她向每个来阅览室看书的孩子收10美分，每天从下午2点30分到4点30分，阅览室都会开放。读者呢，包括邻家的孩子，他们可以在这两个小时内看个够。对于他们来说这是相当便宜的，因为10美分只能买1本连环画，而两小时足够他们看五六本了。

迈克的妹妹在读者离开时要负责检查，确保他们不把书带走。她还要保管书，记录每天有多少人来、他们的名字以及他们有什么意见。我和迈克在以后三个月里平均每星期能赚9.5美元，我们付给迈克的妹妹1美元，而且允许她免费看书，但她很少看，因为她总是在学习。

我和迈克仍然每星期六去小店干活，从各个店收集卖不出去的连环画。我们恪守了对书商的诺言，没有卖一本连环画，如果书太破了我们就烧掉。我们想开一家分店，但实在找不到一个像迈克的妹妹那样尽职尽责、值得信任的管理员。

小小年纪，我们就已经发现找个好职员有多么困难了。

阅览室开张3个月后，发生了一场争斗，附近的小流氓盯上了这桩生意。富爸爸建议我们关门，所以我们的连环画生意就此结束了，同时我们也不去小店工作了。不管怎样，富爸爸十分兴奋，因为他可以教我们新东西了。他很欣慰，因为我们的第一课学得这么好。我们已经学会怎样让钱为我所用了。由于没有从小店的工作中得到报酬，我们就不得不发挥想象力去寻找挣钱的机会。通过经营我们自己的连环画阅览室，我们就掌控了自己的财务，而不是依赖雇主。最棒的是我们的生意让我们赚了钱，甚至当我们不在那儿时，它也照样赚钱，我们的钱为我们工作了。

虽然富爸爸没有付给我们工钱，却给了我们更多的东西。

